\section{Limitations and Ethics}\label{sec:limitations}

\subsection{Coverage limits}

The corpus should not be described as a complete census of all Slovene Bluesky posts ever published. It covers currently recoverable public posts from discovered Slovene-linked accounts; deleted posts are excluded by design. The false-negative analysis (Section~\ref{sec:validation}) estimates an 8.5\% miss rate within the pool of excluded untagged posts from known authors that met the project's minimum-text rule; the true miss rate across the full user population is unknown and likely higher.

The methodology uses post-level rather than thread-level language classification. Slovene-linked users may post in English or other languages within otherwise Slovene conversations, and such replies are not automatically treated as Slovene corpus material.

\subsection{Open data and reproducibility}

The AT Protocol's openness is what made this corpus possible. The collection described here used public protocol mechanisms rather than a paid or application-only research API. This does not mean that access is unlimited: Bluesky documents rate limits for its services, and the ATProto sync specification explicitly allows service providers to apply rate limits and data quotas \parencite{atproto-sync,bsky-ratelimits}. For a small-language community like Slovene, however, the availability of public synchronisation mechanisms is still a practically valuable property, though it also raises specific questions about redistribution.

\subsection{Privacy and identifiability}

The same openness that makes collection easy raises a concern worth stating plainly. The Slovene Bluesky community, at the time of collection, consists of 432 identified authors; a small number. In most cases, the individuals behind the accounts can be identified with little effort: handles are publicly visible, many users link their real names or professional identities through custom domains, and the posting history of any individual is freely accessible through DID resolution. The corpus therefore represents not just a dataset but a detailed longitudinal record of specific, identifiable people.

This is not unique to Bluesky or to this corpus; it is a structural feature of public social-media research. \textit{Public} does not mean \textit{anonymous}: social-media content can remain identifiable even when obvious identifiers are removed, and direct quotation or full-text redistribution can make reverse identification easier \parencite{aoir2019internet,ayers2018dont}. Users who post under their real names may not have anticipated that their posting history would be compiled into a citable research corpus. Following standard practice, full post texts are not reproduced in this article, and any future redistribution should take a conservative form to limit re-identification risk.

\subsection{Corpus redistribution on a decentralised platform}

Whether and how to redistribute this corpus is not straightforward, and this paper does not attempt to resolve the question. ATProto supports content and account deletion, and its synchronisation specification explicitly warns against public redistribution of static repository snapshots in bulk, as mirrors are expected to respect record deletions and account-status changes \parencite{atproto-sync}. A static corpus distributed as a file cannot fully reflect this behaviour: posts deleted on Bluesky remain present in copies that have already been downloaded elsewhere. From a European data protection perspective, public posts may still constitute personal data when authors are identifiable \parencite{gdpr2016}.

Identifier-based redistribution has been used in Twitter research to balance reproducibility with deletion and access constraints \parencite{kinderkurlanda2017archiving}. For ATProto, an equivalent and widely accepted redistribution model has not yet emerged. In the present study, the corpus is therefore described but not released as a full-text dataset. The methodology and collection code are made available, and it is in principle possible to reassemble the corpus from public ATProto data under the conditions described in this paper.

At the same time, it is the author’s intention to make the corpus available in a form that is legally and ethically defensible. A carefully defined redistribution model would enable reuse without requiring individual researchers to repeat the full collection process on their own. However, determining an appropriate access format for a small and identifiable language community requires further consideration and is treated as future work rather than resolved here.
